% Part: lambda-calculus
% Chapter: introduction
% Section: minimization

\documentclass[../../../include/open-logic-section]{subfiles}

\begin{document}

\olfileid[pt]{lam}{int}{min}
\olsection{As funções \usetoken{S}{lambda definable} são fechadas sob minimização}

\begin{lem}
Suponha que $f(x,y)$ seja !!{lambda definable}. Seja $g$ definida por
\[
g(x) \simeq \umin{y}{f(x,y)}.
\]
Então $g$ é !!{lambda definable}.
\end{lem}

\begin{proof}
A ideia, em linhas gerais, é a seguinte. Dado $x$, usaremos o termo lambda
de ponto fixo $Y$ para definir uma função $h_x(n)$ que procura um~$y$ a
partir de~$n$; então $g(x)$ é simplesmente $h_x(0)$. A função~$h_x$ pode
ser expressa como solução de uma equação de ponto fixo:
\[
h_x(n) \simeq
\begin{cases}
n & \text{se $f(x,n) = 0$} \\
h_x(n+1) & \text{caso contrário.}
\end{cases}
\]

Eis os detalhes. Como $f$ é recursiva primitiva, ela é !!{lambda defined}
por algum termo $F$. Lembre-se de que também temos um termo lambda~$D$ tal
que $D(M, N, \bar{0}) \red M$ e $D(M, N, \bar{1}) \red N$. Fixando $x$
por enquanto, para !!{lambda define} $h_x$ queremos encontrar um termo~$H$
(dependente de~$x$) que satisfaça
\[
H(\num{n}) \equiv D(\num{n}, H(S(\num{n})), F(x, \num{n})).
\]
Podemos fazer isso usando o termo de ponto fixo~$Y$. Primeiro, seja $U$ o
termo
\[
\lambd[h][\lambd[z][D(z,(h(Sz)),F(x,z))]],
\]
e então seja $H$ o termo~$YU$. Observe que a única variável livre em~$H$ é
~$x$. Mostremos que $H$ satisfaz a equação acima.

Pela definição de $Y$, temos
\[
H = YU \equiv U(YU) = U(H).
\]
Em particular, para cada número natural~$n$, temos
\begin{align*}
H(\num{n}) & \equiv U(H, \num{n}) \\
& \red D(\num{n}, H(S(\num{n})), F(x, \num{n})),
\end{align*}
como requerido. Observe que, se você tomar um numeral $\num{m}$ como substituto de~$x$
na última linha, a expressão se reduz a $\num{n}$ se
$F(\num{m}, \num{n})$ se reduz a $\num{0}$, e se reduz a
$H(S(\num{n}))$ se $F(\num{m}, \num{n})$ se reduz a qualquer outro
numeral.

Para concluir a prova, seja $G$ o termo $\lambd[x][H(\num 0)]$. Então $G$
!!{lambda define}s~$g$; em outras palavras, para todo~$m$, $G(\num m)$ se
reduz a~$\overline {g(m)}$, se $g(m)$ está definida, e não tem forma normal
caso contrário.
\end{proof}

\end{document}
